\paragraph{Why it stops halfway.} The reason is baked into how the operator picks winners. It
runs a separate max \emph{per filter}: each filter looks at its own activations and keeps its
strongest spot in each window. Crucially, the filters never talk to each other. Nothing in the
math says ``filter B, that position is taken --- go somewhere else.'' So when a sequence
position is genuinely informative, every filter that responds to it keeps it, and they all pile
on. The operator thins out each filter's \emph{own} hits (within-filter non-overlap, which it
guarantees), but it never divides the territory \emph{between} filters.

\paragraph{Two levers that help --- and one that would finish the job.} Experiment~3 shows the
softmax strength $\alpha$ is a \emph{soft} version of exactly the fix we need. Raising $\alpha$
sharpens the over-channels softmax toward a per-position winner-take-all, and it does move
filters apart: occupancy falls from $2.8$ toward $2.4$ and exclusivity nearly triples. There is
a sweet spot around $\alpha=20$ where this separation is free of any held-out accuracy cost;
only at $\alpha=100$ does accuracy start to slip. A second, independent lever is simply
\emph{fewer filters}: dropping the inference layer from 20 to 10 filters (the bottleneck)
halved the crowding on its own. But neither closes the gap --- even $\alpha=100$ leaves
occupancy at twice chance. The \emph{hard} version would: a single winner per position across
all filters (a max \emph{across} channels instead of within each), a one-line change that forces
the territory to be split. Its risk is over-separation --- it may hurt tasks where nearby
filters \emph{should} co-fire, such as discovering motifs that genuinely co-occur at one site.

\paragraph{Annealing $\alpha$ is the natural next step.} The one place high $\alpha$ hurt was
training from scratch at $\alpha=100$. Ramping $\alpha$ up \emph{during} training --- let the
model first learn to fit at $\alpha=1$, then sharpen --- should reach high-$\alpha$ separation
without paying the accuracy tax. The knob is already in place (\texttt{sparse\_unpool\_alpha});
only a per-epoch schedule is missing.

\paragraph{The downstream cost is real.} At full length (Experiment~2), the sparsification
\emph{hurt} motif discovery: triplet groups collapsed from $11$ to $0$ and significant pairs
slipped from $8$ to $7$. This is the flip side of separation. The method finds co-occurring
motifs by looking for filters that fire together; the operator's whole job is to stop them
piling up, so it erases the joint firing the higher-order groups need. Separation and
co-occurrence discovery are, on this dataset, in direct tension --- which reframes the goal:
if the aim is co-occurrence, mild separation (small $\alpha$, or none) is safer than pushing it
hard.

\paragraph{Caveats.} Single dataset, single seed, an MSE regression proxy for training (not the
method's full objective), and one inference layer ($L^\star=1$, $p=5$). The occupancy gap
relative to chance is large and consistent across reruns, so the qualitative verdict is robust;
the exact numbers are not.
